\section{承受的边界：后果回到哪里}

行动一旦发生，后果就要落到具体的人、组织和环境中。产生后果的位置与承受后果的位置可能分离，控制权与责任也可能被分散。本节所讨论的责任主体，需要有稳定身份、可追索资源、记录和后果回流机制。人的人格地位与基本权利另有伦理和制度依据，责任结构则回答一次行动由谁解释、纠正和赔偿。

\subsection{主体的边界：后果回流与责任显现}

行动把残余不确定性推入现实之后，承付必须落在一个边界上，而“边界”本身在思想史上有很长的来路。在人类很长一段思想史里，最先被追问的不是主体，而是世界本身。什么存在，什么是真实，变化背后有没有稳定结构，万物如何成为万物，这是本体论的问题。那时，人当然也在世界之中，但人还不是问题的中心；人首先面对的是一个比自己更大的秩序：自然、神、实体、形式、规律、宇宙结构。问题的重心在“世界如何成立”。

\subhead{从本体论到主体论}

后来，问题开始转向人自身。世界不只是存在在那里，还要被经验、被认识、被判断、被表达。于是，主体被推到前台。不是因为世界消失了，而是因为世界进入人的经验和判断以后，必须追问：谁在认识？认识凭什么可靠？判断从哪里获得确定性？这个转向把人从世界中的一个位置，变成了世界显现和意义生成的中心之一。主体不再只是存在者之一，而是世界如何被把握的条件\cite{kant1781kritik}。

再往后，问题又从认识推进到行动。人不只是看世界、解释世界，也在改变世界。技术、工业、科学实验、组织、工程和制度，把人的判断变成了现实中的力量。于是，仅仅问“谁在认识”已经不够，还要问“谁在行动”。行动不同于认识，因为行动会改变状态，会制造不可逆后果，会让世界不再只是对象，而成为被干预、被重构、被损伤或被修复的过程。现代智能把这个问题继续往前推了一步。智能系统不只生成表述，不只辅助认识，也开始参与行动：它可以生成判断，触发流程，操控机器，参与实验，影响交通、医疗、金融、教育和生命过程。问题因此不再只是“谁在认识世界”，也不只是“谁在改变世界”，而是：改变发生以后，后果回到哪里？

这就是从本体论到主体论，再到行动论和承受论的真正来路。本体论问世界如何存在，认识论问世界如何被把握，行动论问世界如何被改变，而主体论在智能时代必须追问：改变之后，反作用由谁承受。过去，意图、身体、行动和后果大多绑定在同一个生命或制度边界里；现在，生成、执行、收益和损害可以分散到不同位置。行动发生了，世界被改写了，反作用却未必回到制造差异的地方。所以，主体问题不是因为机器开始像人说话才重新尖锐，而是因为智能进入现实以后，作用与反作用开始脱钩。一个边界如果只向世界施加差异，却不被差异的后果反过来改写，它就还不是完整主体，而只是无主行动中的一个发生点。\key{主体不是意识叙事的中心，而是后果回流的边界。}

上一节讨论的是行动准入：一个候选、判断、方案或指令，什么时候可以跨过表示层，进入现实系统。这里追问的是下一步：一旦行动进入现实，反作用回到哪里？没有准入，行动会过早发生；没有主体边界，行动会无人承受。智能时代的问题不是行动太少，而是行动越来越容易被生成、触发和放大。真正稀缺的，是让行动重新接回承受边界的能力。

\subhead{后果回流的边界}

主体不是会说“我”的东西。会自称，会解释，会表达偏好，会生成计划，会承认错误，都只能制造主体的外形。主体也不是单纯能行动的东西。能行动只说明它能向世界施加差异，不能说明差异造成的反作用会回到它自身。

主体首先是作用与反作用闭合的地方：行动从这里发出，后果回到这里，代价在这里沉积，未来行动因此被约束。所以，主体不是行动的起点，而是后果的归宿。一个大模型可以输出判断，一个具身系统可以执行动作，一个生命智能可以提出干预方案；这些都还不够。关键在于，错误、损伤、失败、羞耻、信任破裂、声誉变化、身体代价、时间损失和未来限制，是否会回到同一个持续边界中。没有后果回流，就只有输出；没有代价沉积，就没有主体。

这里必须区分反馈和承受。反馈可以很轻。温控器有反馈，控制系统有反馈，模型可以根据错误更新参数，机器人可以根据碰撞调整路径\cite{wiener1948}。反馈回答的是系统如何修正，如何提高效果。承受不同。承受意味着代价回到行动者身上，留下不可随意擦除的痕迹，并改变它未来能够如何行动。参数更新不是悔恨，路径重规划不是负责，损失函数下降不是付出代价。反馈是信息回流，承受是代价回流。主体从承受开始，而不是从反馈开始。

承受从哪里开始，人的身体给出了最直接的例子。人的身体之所以天然形成主体边界，不是因为人有某种神秘特权，而是因为人的生命过程足够难以逃离\cite{merleauponty1945}。一个人犯错，会痛，会羞耻，会失去时间，会破坏关系，会改变声誉，会在记忆中反复返回那个后果。人不能像模型一样轻易重置版本，不能像零件一样无痛替换，不能把衰老、伤害和死亡外包给另一个实例。人之所以天然接近主体，不是因为人更高贵，而是因为人不能无成本逃离自己。

\subhead{两种承受封装}

所谓\term{承受封装}，就是后果不能随意散开。身体记得伤口，神经记得恐惧，习惯记得重复，关系记得背叛，声誉记得选择，时间记得不可逆。一个“我”不是一句自我叙述，而是后果在身体、记忆、关系和时间中不断回流、折叠、压实之后形成的持续边界。主体不是静止内核，而是被过去行动不断改写的连续性。生命主体不能每次行动后彻底清零，只能带着沉积过的后果继续行动。

具身智能容易制造一个误解：机器有身体、能感知、能行动，所以就获得了主体性。这不成立。机器人的外壳、关节、传感器和执行器构成物理边界，但物理边界不自动等于主体边界。机器可以损坏，零件可以更换，模型可以替换，机体可以报废，数据可以回收，成本可以转移。许多后果并不回到机器自身，而是散落到制造者、使用者、部署者、环境和被影响的人身上。身体如果只是执行器，就不是主体边界；身体只有成为后果沉积的地方，才具有主体意义\cite{dreyfus1972what}。

生命体不同。生命体的身体不是临时接口，而是长期演化形成的承受结构。许多复杂性被封装在形态、反射、疼痛、恢复、习惯和有限寿命中。机器没有自然继承这笔演化摊销。它可以拥有行动能力，却不必然拥有承受能力；它可以进入空间，却未必把空间中的反作用带回自身；它可以改变身体，却未必承受身体被改变后的余波。这就是具身性和主体性之间必须划开的边界。

反过来，没有自然身体，也不意味着完全不能形成主体。一个组织没有疼痛，没有肉身，没有生命意义上的死亡，但它可以通过资产、账目、历史记录、声誉、身份和长期约束形成某种制度边界\cite{dewey1926historic}。它过去的行动会改变它未来能做什么，能获得什么信任，能承担什么限制。后果可以在制度中沉积，制度也可以因此被过去改写。这说明主体不是生物身体的专利。主体可以由生命封装，也可以由制度封装。关键不在材料，而在后果能否在一个边界内持续闭合。

\subhead{无主行动与责任显现}

闭合成立与否，在多数智能系统上一眼可见。当前多数智能系统的问题，不是它们“不像人”，而是它们的闭合边界尚未成立。大模型能生成判断，但判断造成的后果不自然回到模型；具身系统能执行动作，但动作的反作用不一定沉积在机器自身；生命智能能提出干预方案，但被分类、筛选、治疗和优化的人，才真正承受长期结果。智能越能生成、行动和干预生命，越会放大这个缺口：差异由系统制造，反作用却由别处承受。这不是简单的“谁负责”问题，而是更底层的“行动是否有主”问题。没有主体边界，行动会退化为输出；没有承受闭合，现实会退化为反馈数据；没有代价回流，错误会被分摊到看不见的地方。一个系统可以越来越有效，却越来越不承受；可以越来越会行动，却越来越像没有主人。要防止的，正是这种世界被纯输出化的趋势。

责任显现也不应被理解成外部条款。它的深层意义，是让被系统吞掉的后果重新可见。当行动只影响行动者自身，反作用常常自然显现：摔倒会疼，失败会记住，关系会破裂，时间无法倒流。但当行动通过模型、机器、组织和生命技术放大到更大范围，后果就可能被延迟、遮蔽、统计化、转嫁，变成没有一处真正疼痛的损失\cite{matthias2004responsibility}。责任不是事后装上的道德标签，而是主体边界在社会中的显现。它让行动不能只留下效果，也必须留下归处。

主体性因此不是非有即无的标签，而是后果回流的闭合强度。工具几乎没有闭合，锤子伤人，后果回到使用者。普通代理体有目标和反馈，却未必有承受。制度主体没有自然身体，却可以让后果在组织记忆和长期限制中沉积。生命主体拥有最强的承受封装，身体、记忆、关系、衰老和死亡把后果锁在同一个生命过程中。未来也可能出现工程化主体，但那不取决于它会不会说“我”，而取决于它是否拥有持续身份、不可任意重置的记忆、真实代价沉积、外部可追问边界和长期未来约束\cite{parfit1984reasons}。主体论不是为了捍卫人的神秘性，也不是为了预先否定机器可能获得某种主体地位。它要提出的是判准：谁被行动的反作用持续改写，谁才开始接近主体；谁只负责施加差异，却总能把代价转走，谁就仍然只是通道。智能时代真正危险的不是机器像不像人，而是行动越来越多，后果越来越散，世界开始出现大量无主行动。

主体不是会说“我”的东西。主体是后果不能再继续转嫁的地方。它把行动、反作用、记忆、代价和未来约束重新扣在一起。没有这个边界，智能系统会越来越能改变世界，却越来越难被世界反过来改变；行动会越来越多，主人却越来越少。但还有一个入口没有被检验。机器可以拥有身体，能感知、能执行、会磨损、会自我维护，这些能力常被读作主体性的证据；如果主体性只取决于后果能否在一个边界内闭合，那么材料边界与承受边界就必须分开来看，而工程上完全可以把身体做强的那些条件，恰恰未必等于承受闭合形成的条件。

\subsection{承受闭合：从机器身体到工程化主体}

而工程上出现了一个更容易误判的问题：机器有了身体，能感知、能行动、会磨损、会自我修复，是不是就更接近主体。这个问题之所以重要，是因为智能进入现实有几种不同入口。大模型让输出进入社会判断，具身智能让身体进入材料世界，生命智能让方案进入代谢、发育、免疫、遗传和生态过程。它们都比单纯的屏幕文字更重，也都更容易让人误以为：智能越接近现实，就越接近主体。

\subhead{现实入口不是主体边界}

但现实入口不是主体边界。一个系统可以输出判断，却不承受判断造成的后果；可以拥有身体，却不承受身体行动造成的代价；可以提出生命干预方案，却不承受代谢、发育和代际后果。进入现实只说明作用发生了，不说明反作用已经回到同一个持续边界。主体性不从输出中自然长出，不从身体中自然长出，也不从生命干预能力中自然长出。主体性只在一个地方显现：后果是否回流、沉积，并约束同一边界的未来。

因此，首先要把身体、具身性和主体性拆开。身体是材料边界，回答系统以什么物理结构进入世界。具身性是材料边界与环境形成行动耦合，回答系统如何感知、耗能、受扰、修正和继续行动。主体性是后果回流和承受闭合，回答行动之后，代价是否回到一个不能随意逃离的连续边界。三者相连，但不能互相替代。

身体可以被工程化。传感器可以增加，执行器可以升级，关节可以维护，材料可以柔顺化，能量可以调度，局部控制可以低延迟，自诊断可以常态运行。未来机器身体甚至可能具备更强的自维护和自修复能力：检测磨损，修补表面裂纹，重新分配能量，在局部损伤后降级运行\cite{terryn2017self}。但这些首先仍然是身体工程，不是主体生成。

上述区分可以写成一个只刻画回流覆盖面的比率。设 $B$ 是一个持续边界，$a$ 是它发起的行动，$C(a)$ 是该行动的全部后果，$\rho_a$ 把每一项后果映射到最终承载它的边界。定义 $B$ 在行动 $a$ 上的后果回流比例

\begin{equation}
\kappa(B,a)=\frac{\bigl|\{\,c\in C(a):\ \rho_a(c)=B\,\}\bigr|}{\lvert C(a)\rvert},
\qquad \kappa\in[0,1].
\label{eq:ch33-closure}
\end{equation}

这里假定$C(a)$是有限、非空且分项一致的集合，$\rho_a$把每项归于一个边界。$\kappa$只记录这份清单上的回流覆盖。实际后果可以共同承担，也可能无法计数；分项方式和权重都会改变结果。这个比率帮助追问归属，不测量意识、道德地位或完整主体性。保险和公共救济还可能正当地分担后果，不能把全部回流当成普遍目标。

\subhead{残差吸收与后果承受}

这里必须区分残差吸收和后果承受。具身系统进入物理世界以后，滑移、摩擦、惯性、接触误差、传感漂移和任务链偏差都必须被处理。它们可以被身体吸收，被控制器修正，被环境约束，被外部支撑结构分担。残差吸收解决的是行动能否继续。后果承受追问的是行动归属于谁。具身性容易被误读，正因为它把这两个问题压得很近：机器真的在处理现实残差，于是看起来也像在承受现实后果。

但一个东西被现实改变，不等于它在主体意义上承受现实。石头会风化，桥梁会疲劳，轮胎会磨损，电池会衰减，数据库会被写入，模型会被微调，机器会留下日志。它们都在时间中留下痕迹，但这些痕迹未必构成一个行动者的历史。被改变只是物理事实，承受则意味着后果被纳入一个可追问的连续性。这里的关键差异，是承载和承受。轮胎承载摩擦，机械臂承载反力，电池承载放电循环，外壳承载碰撞和磨损。只要机器进入物理世界，后果就会落到某些部件上。但承载只是后果落在\term{承载面}上；承受要求这些后果进入同一持续边界，改变它的未来可达空间，并使它不能无成本逃离过去。

一个服务机器人撞坏墙角，外壳可能划伤，轮组可能偏移，日志可能记录异常，控制策略可能更新。物理后果落在了机器上，信息反馈也回到了系统里。但真正的反作用可能落在用户、制造者、部署者、空间环境和被影响的人身上。机器身体只是后果链中的承载面，不必然是后果链的归宿。它受到了影响，却未必成为承受者。所以，承受至少包含几项条件：后果回到同一边界，而不是只落到可替换部件；后果改变这个边界的未来可达空间，而不是只触发维修或参数调整；行动带来的增益、限制和反作用约束同一连续体，而不是向不同位置散开；这个边界能够被外部持续追问，而不是在版本、实例、供应链和维护链条之间消失。少了这些条件，机器身体即使有磨损，也只是后果经过的材料界面，不是后果归属的主体边界。

\subhead{自我保护与生命封装}

很多讨论把“会保护自己”看成主体性的证据，这同样需要拆开。机器人会避障，不等于它害怕受伤；会低电量返航，不等于它关心未来；会检测故障，不等于它理解自己的脆弱；会自我修复，不等于它承受了创伤。避障可以是控制目标，返航可以是能量策略，自修复可以是材料机制。它们可以很复杂，也可以很有效，但它们首先是在维持功能，不是在形成一个不能逃离自身的“我”\cite{metzinger2003being}。生命体为什么不同？不是因为生命体内部藏着一团神秘主体，而是因为生命把承受封装得更深。

身体、代谢、疼痛、恢复、发育、记忆、关系和死亡，把后果压进同一个不可随意替换的过程。身体受伤不是单个部件异常，而会改变行动能力和未来时间；疼痛不是报警灯，而是把损伤带入同一个生命边界；死亡不是报废外壳，而是连续性的终止。生命不是更高贵，而是更难把代价转走。生命体不能随意更换完整身体，不能把记忆无痛迁移到另一个实例，不能在失败后从干净版本重新启动，不能让一个等价副本替自己承受衰老和死亡。主体性不是从会动的肉体里长出来，而是从不能无成本逃离自身的连续性中长出来。

机器身体没有自然继承这条连续性。它可以有编号、位置、外壳、传感器和维护记录，却仍然可能只是一个可替换实例。关节磨损了，可以换；模型出错了，可以回滚；机体报废了，可以用同款机器接替；策略失败了，可以在新版本中擦除。机器当然可以有历史，问题在于，这些历史常常属于平台、车队、制造商、外部运行结构或数据库，而不属于执行个体自身。历史被记录，不等于历史被承受。

\subhead{承受闭合的条件}

这不是说机器永远不可能形成主体性，而是说主体性不能由身体自动推出。未来若出现某种工程化主体，它需要的不是更像人的外形，也不只是更精密的触觉皮肤，而是承受闭合的工程：持续身份，不可任意重置的历史，真实代价沉积，未来行动受过去后果约束，外部能够追问的边界，以及不能把所有代价转嫁给用户、部署者、环境和受影响者的结构。这样的结构可以不是生物性的。组织没有自然身体，却可以通过资产、声誉、法律身份、历史记录和长期约束形成制度边界。它没有疼痛，却有后果沉积；没有肉身，却有身份终止。主体性不是肉体专利，但它要求后果不能随意散掉。

这也解释了智能进入现实的三种错位。大模型的错位，是表示能够影响现实判断，却不承受判断后果；具身智能的错位，是身体能够承接现实残差，却不必然归属行动代价；生命智能的错位，是方案能够进入生命过程，却未必承受代谢、发育、遗传和生态的长后果\cite{floridi2013distributed}\cite{nissenbaum1996accountability}。越往后，承受面越重，系统自身却未必更像主体。风险不在于智能越来越像人，而在于智能越来越能进入高后果层，却仍然可以把代价留在别处。

因此，判断一个智能系统是否接近主体，不应看它像不像人，也不应只看它有没有身体，而要看它是否形成了承受闭合。一个系统可以高输出而低主体，可以高度具身而低主体，也可以深入生命过程而仍然没有承受闭合；反过来，一个组织没有自然身体，却可以通过制度形成主体边界。表示是现实影响的入口，身体是材料耦合的入口，生命技术是生命过程的入口，主体性则是后果回流和承受连续的闭合强度。

\subhead{耦合、自我修正与可替换性}

更底层地看，智能系统进入现实有几种不同耦合。表示耦合回答系统能生成什么候选、影响哪些判断。材料耦合回答系统能否在物理世界中稳定行动。生命耦合回答一个方案进入生命过程以后，会被什么生命结构读取、放大、抵抗、沉积或传播。承受耦合则回答后果是否回到一个持续边界。前三种耦合说明智能如何进入现实，最后一种耦合才触及主体性。

这也解释了为什么“自我修正”仍然不是主体性。大模型可以反思输出，具身系统可以监测磨损并主动安排维修，生命智能系统可以根据实验结果更新候选、优化剂量、调整干预路径。它们看起来都在改正自己，但改正不等于承受。只有当错误不能只是变成新版本、新策略、新疗程或新候选，而会以不可转嫁的方式改变同一持续边界的未来，修正才开始接近承受。

更深的分界在于可替换性。身体作为工程对象，天然倾向于模块化、可维护、可替换、可复制。工程追求降低单点不可替代性：坏了能换，错了能回滚，实例能迁移，能力能复用，影响能隔离。这对可靠性是好事，却恰恰削弱了主体性的自然生成。主体需要某种不可完全替换的连续性，而工程常常通过可替换性提高系统韧性。机器越是良好的工程对象，越未必是主体。

这不是矛盾，而是两种目标的分叉。工程身体追求可维护，主体边界要求可追问；工程系统追求故障隔离，主体性要求代价沉积；工程结构追求实例迁移，主体性要求身份连续；工程优化追求外部影响最小化，主体判准追问反作用最终回到哪里。身体可以被设计得越来越强、越来越灵活、越来越自洽，但只要后果仍然由别处承受，它就还是进入现实的通道，而不是后果的归宿。

\subhead{不可逃离的边界}

这正是“身体可以工程化，主体不能自然继承”的精确含义。工程化身体的目标，是把现实扰动处理掉，让系统少停机、少损坏、少失控；主体边界的目标，不是把后果处理掉，而是让后果不能失踪。一个方向追求残差被吸收，另一个方向追问代价由谁保留。工程越成熟，越可能把故障消隐在维护、替换、冗余和外部支持结构中；主体性越清晰，越要求这些被消隐的后果重新回到可追问边界。两者可以相互支持，但绝不能互相推出。

最硬的判准在这里：\key{主体性不看系统能否保护自己的部件，而看它能否逃离自己的后果。}一个系统如果能把伤害分散给用户，把风险交给外部结构，把错误归入版本，把损坏交给维护，把身份迁移到新实例，而自己继续作为干净能力出现，它就还没有成为主体。相反，如果某个边界不能这样逃离，过去行动会持续限制它、定义它、追问它，并改变它未来能够进入的路径，那么即使它不是生物，也已经开始接近主体。

主体不是身体的自然升级版。身体让机器进入世界，具身性让机器在世界中形成行动闭环，生命技术让智能进入更重的生命过程；但主体性要求世界把行动的后果带回同一个持续边界。它不从传感器里长出，不从执行器里长出，不从生命干预能力里长出，也不从自我描述里长出，而从代价无法继续转嫁的地方显现。因此，主体性不能从任何一种现实入口中自然推出。输出、身体、反馈、自维护和生命干预，都只能说明系统已经以某种方式接入现实；它们还没有说明后果会回到哪里。只有当反作用持续回到同一个边界，并改变它未来能够如何存在、如何行动、如何被追问时，这个边界才开始具有主体意义。具身性让机器进入现实，承受性才让一个边界成为主体。

不过，单个边界的闭合只是判准的一半。一次行动很少只穿过一个边界：平台、开发者、部署者、组织、用户和受影响者各自承担了一部分后果，回流路径不止一条，收束点也不止一个。当后果分散在多条回流路径上时，闭合究竟在何处成立，谁来决定哪个边界承担哪一部分，代价又如何在多个承受者之间分配，这些还没有答案。

\subsection{分布式承受：为什么公司能负责而模型不能}

上一节停在了一个没有答案的问题上：当后果分散在多条回流路径上时，闭合在何处成立，谁来决定哪个边界承担哪一部分。要回答它，可以先看一个已经把这件事做完的例子。有一种主体，不吃饭，不呼吸，没有神经系统，没有经历，也没有痛苦。它可以拥有财产，签订合同，被起诉，被罚款，被强制解散。它要为自己的行为负责。这种主体叫公司。

\subhead{一个已经解决的难题}

公司承担责任这件事，今天已经不再让人惊讶，以至于它的构造性被忽略了。公司不是自然物。它是一组法律装置人为组装出来的承受边界。人们花了几个世纪才把这组装置拼齐，其中最关键的一步是有限责任：投资人以其出资额为限承担责任，公司以其全部资产承担责任，两者分开。有限责任常被当成一种鼓励投资的经济政策。它在结构上的作用更基础：它把一条本来会无限延伸的后果链，截断在一个可以清点的边界上。在这条边界内，有资产可以扣，有账户可以查，有代表可以传唤，有主体可以解散。因为边界存在，追责才成为一个可以操作的动作，而不是一场无限追溯。

这提供了一把尺子。用它可以量出，一个系统要成为后果的承受者，需要什么条件。用这把尺子去量机器智能系统，缺什么就很清楚了，但在那之前，先要把这把尺子本身说清楚。

\subhead{承受边界的四个条件}

把公司的结构拆开，可以看到四件事同时在起作用。第一是\term{可清点的资源边界}。公司有自己的资产和负债。赔偿不是象征性表态，而是从资产负债表上真实划走一笔。资源边界让惩罚具有量纲。第二是\term{稳定的同一性}。公司持续存在，跨年度、跨管理层、跨员工流动。今天的公司要为十年前的公司行为负责，因为它们被认定为同一个法律人格。同一性是追责的时间通道。

第三是\term{受罚后的内部状态改变}。罚款会减少可投资金，影响信用评级，抬高融资成本，压缩薪酬空间，甚至触发破产。也就是说，惩罚会真正改变这个边界未来能够做什么。这是承受的实质：不是被指向，而是被改变。第四是\term{可追索的责任链}。公司有董事、高管、法人代表、雇员和股东，各自有明确的职权范围和责任分配。

一次行为可以被拆解到具体的决策位置，而不是消失在集体里。这四条加在一起，构成了一个完整的承受结构。\key{公司能负责，不是因为它有意识，而是因为制度把后果固定到了一个可追索边界上。}这四条不是同等重要的四条，而是缺一不可的四条：任何一条空着，承受结构就不成立。

\subhead{模型缺什么}

用同一把尺子量机器智能系统，缺什么就很清楚了。第一，缺资源边界。一个模型没有资产。它可以被罚款吗，罚款由谁支付？如果由部署它的公司支付，那承受者就是公司，模型本身没有承受任何东西。这一点看似只是会计问题，实际上是结构问题：没有可清点的资源，惩罚就无法落到系统身上，只能落到系统之外的人身上。

第二，要查明同一性由谁保存。模型版本变化不必然切断责任：部署者可以保留版本、配置、记录与持续义务。真正的问题是改版以后，既有行动能否仍被定位，承担义务的边界是否继续存在。

第三，要查明制裁实际改变谁。罚款减少机构资源，服务限制改变运营能力，模型更新则改变行为。几种反应可以相互连接，但承担费用和义务的通常是人或机构。把它们统称为“模型受罚”，会遮住实际落点。

第四，缺可追索的责任链。这不完全是技术问题。模型的生成、微调、部署、集成、调用和使用，常常分布在不同的主体之间。一个输出进入现实，中间可能经过模型提供方、应用开发者、平台运营方、终端使用者和受影响的第三方。每一段都有责任缺口，而责任在传输中会衰减。

把这种衰减写成逐层传递的形式，

\begin{equation}
R_n=R_0\prod_{i=1}^{n}\alpha_i,\qquad 0\le \alpha_i\le 1
\end{equation}

其中 $n$ 是行动经过的层数，无量纲，正文给出的典型情形是六到七层；$R_0$ 是行动刚发生时可归属到行动位置的责任，为便于比较可归一化为一；$\alpha_i$ 是第 $i$ 层向下一层传递之后、仍然能够归属到具体决策位置的责任份额，无量纲，取值在 0 与 1 之间；$R_n$ 是经过全部层之后仍然落在某个可追索边界上的责任份额。这个式子的含义是：只要每一层都截留一部分，剩余份额就随层数按乘积下降，层数越多下降越快，这正是“责任在它们之间分散、转移、稀释，最后变成无人完整承受的残差”的写法。它也解释了集中归属与举证责任倒置为什么有效：这两种制度处理并不减少层数，而是从外部把某些 $\alpha_i$ 抬回到接近 1。适用条件是各层贡献可分离，且不在多层之间重复归责；当同一主体同时提供模型、应用与平台时，乘积链条合并，$n$ 减小，衰减随之减弱。它标出的是结构位置，不是实测比例，$\alpha_i$ 本身无法直接测量。

通常部署的模型是责任安排中的一个组件，而非独立的法律承受者。持续身份、资源与权限可以由外部架构配置，义务和救济则需由制度确定。问题应逐项落到具体系统：谁控制、谁记录、谁获益、谁补偿。不能从当前模型缺少某种安排，推出任何工程化边界原则上都无法具备它。

\subhead{分布式承受}

真正的问题因此不是“模型能不能负责”，而是当行动分散到多个位置之后，后果应该被固定到哪里。这个问题在几个领域已经出现过。自动驾驶的碰撞责任要在整车制造商、软件提供方、传感器供应商、高精地图服务商、车主和驾驶员之间分配。同一场事故，不同的技术架构会导出不同的责任划分。

内容平台的责任要在发布者、平台、推荐算法和用户之间分配。平台以中立通道自居时，推荐算法的编排作用如何计入，长期没有一致答案。算法决策造成的差别对待，受害者要证明什么才能立案，被告要证明什么才能免责，举证责任的分配直接决定了这类案件是否可诉。这些案例的共同结构是：行动被拆散了，承受边界却没有同步拆散，也没有重新组装。发起、执行、获益和受损分处不同位置，而后果需要一个落点。

在传统结构里，这个落点由行动者本人担任。一个人做出行动，身体、身份、声誉、时间和未来都押在这条路径上，失败会改变他自己。机器智能接入之后，路径照样可以被启动、被执行、被放大，但这些动作不落在任何一个会因此改变的边界上：模型这一侧没有与行动者本人相当的位置，也没有哪个部分会随失败而改变。缺失的部分不会自动长出来。它只能由制度装置地从外部装上。

而一旦装上，人的位置就跟着变了。分布式承受被解决的过程，同时是责任从个体身上被移走的过程。当系统能够承接后果，个体与后果之间的直接联系就松了一环。松掉的那一环，原本是许多人的价值来源。如果价值曾经附着在“后果落在自己身上”这个位置上，那么松掉的一环就需要新的解释：当制度能够承接后果时，还有哪些判断必须由人来作出。价值不会消失，但它需要一个新的源点。

不过，制度接住后果的过程，最先改变的不是价值的论证，而是每个人对自己位置的判断。在价值需要一个新源点之前，先要发生的是一件更直接的事：许多人发现，自己原来用来证明价值的方式失效了。
